%=====================================================================
\section{The Unique Substrate: Why $\mathbb{C}$}
\label{ch:whyC}
%=====================================================================

\subsection{The starting point}

We begin with the minimal possible assumption about reality:

\begin{quote}
\textbf{Assumption.} \emph{There exist $N$ points. Between every ordered pair $(i, j)$, there exists a relation $G_{ij}$.}
\end{quote}

This is not a physical axiom --- it is the minimal structure needed to speak of ``things'' and ``relationships between things.'' A universe without relations between its constituents is not a universe; it is a collection of isolated points with no structure, no information, and no physics.

The relation $G_{ij}$ takes values in some algebraic structure $\mathbb{K}$. We do not assume what $\mathbb{K}$ is. Instead, we derive it from the requirements that the relations describe a universe containing both \emph{geometry} (the shape of space) and \emph{forces} (interactions between matter).

\subsection{Four necessary conditions}

For the relation matrix $G = (G_{ij})$ to encode a universe with geometry and forces, $\mathbb{K}$ must satisfy:

\begin{enumerate}[label=\textbf{R\arabic*.}, leftmargin=2cm]
\item \textbf{Magnitude} (Geometry). $G_{ij}$ must have a well-defined absolute value $|G_{ij}|$ encoding the ``distance'' or ``similarity'' between points $i$ and $j$. This requires $\mathbb{K}$ to be a \emph{normed} algebra.

\item \textbf{Phase} (Gauge forces). $G_{ij}$ must have a continuous phase $\arg(G_{ij})$ encoding the gauge connection between $i$ and $j$. A universe with magnitude only has geometry but no forces --- no electromagnetism, no nuclear interactions, nothing to make matter interact. This requires the unit-norm elements of $\mathbb{K}$ to form a \emph{connected Lie group}.

\item \textbf{Determinant} (Gravity). For any subset of points forming a simplex, the ``volume'' $\det(G_h)$ must be well-defined as a single real number. This volume becomes the hinge area in Regge calculus, encoding spacetime curvature. This requires $\mathbb{K}$ to be \emph{commutative}: $G_{ij} G_{kl} = G_{kl} G_{ij}$.

\item \textbf{Winding} (Charge quantization). The holonomy $\Phi = \arg(G_{ij} G_{jk} G_{ki})$ around a closed loop must admit discrete winding numbers, so that charges are quantized (electrons have charge exactly $-1$, not $-0.997$). This requires $\pi_1(\text{phase group}) \cong \mathbb{Z}$.
\end{enumerate}

\begin{remark}
These are not four independent axioms. They are four \emph{necessary conditions} for the existence of a universe with geometry, forces, gravity, and discrete particles. A ``universe'' lacking any of these is either structureless (no forces), volumeless (no gravity), or continuous (no particles). We derive $\mathbb{K}$ from these conditions; the conditions themselves are consequences of asking ``what is the minimum structure that constitutes a universe?''
\end{remark}

\subsection{Classification of candidates}

The candidates for $\mathbb{K}$ are constrained by classical theorems of algebra.

\begin{theorem}[Frobenius, 1877 \cite{Frobenius}]
The only finite-dimensional associative division algebras over $\mathbb{R}$ are $\mathbb{R}$, $\mathbb{C}$, and $\mathbb{H}$ (quaternions).
\end{theorem}

\begin{remark}
The octonions $\mathbb{O}$ are excluded by non-associativity. The relation matrix $G = \Psi\Psi^\dagger$ requires associative multiplication; without it, the matrix product $(\Psi\Psi^\dagger)_{ij} = \sum_a \psi^*_{i,a}\,\psi_{j,a}$ depends on the order of operations and is not well-defined.
\end{remark}

We now eliminate $\mathbb{R}$ and $\mathbb{H}$ by the remaining requirements.

\subsection{Proof I: Algebraic elimination}

\begin{theorem}[Uniqueness of $\mathbb{C}$ --- algebraic proof]
\label{thm:C_algebra}
$\mathbb{C}$ is the unique associative division algebra over $\mathbb{R}$ satisfying R1--R3.
\end{theorem}

\begin{proof}
By Frobenius, the candidates are $\{\mathbb{R}, \mathbb{C}, \mathbb{H}\}$.

\medskip
\noindent\textbf{Elimination of $\mathbb{H}$ (R3: commutativity).}
The quaternions are non-commutative: for the basis elements $\mathbf{i}, \mathbf{j} \in \mathbb{H}$, $\mathbf{i}\mathbf{j} = \mathbf{k} \neq -\mathbf{k} = \mathbf{j}\mathbf{i}$. The determinant of a matrix over $\mathbb{H}$ is therefore ambiguous:
\begin{equation}
\det\begin{pmatrix} a & b \\ c & d \end{pmatrix} \stackrel{?}{=} ad - bc \quad\text{or}\quad da - cb \quad\text{or}\quad ad - cb \quad\text{or}\quad \ldots
\end{equation}
These expressions differ over $\mathbb{H}$. The Dieudonn\'e determinant resolves this ambiguity but maps to $\mathbb{R}_{\geq 0}/\mathbb{H}^\times$, losing the sign structure essential for deficit angles in Regge calculus ($\delta_h$ can be positive or negative). Without a well-defined signed determinant, spacetime curvature cannot be encoded. $\mathbb{H}$ is eliminated.

\medskip
\noindent\textbf{Elimination of $\mathbb{R}$ (R2: continuous phase).}
The unit-norm elements of $\mathbb{R}$ are $\mathbb{R}^\times_1 = \{+1, -1\} \cong \mathbb{Z}_2$. This is a discrete group with two elements. A gauge connection requires a \emph{continuous} phase: the gauge field $A_{ij} = \arg(G_{ij})$ must vary smoothly between neighboring points to define a field strength $F = dA$. Over $\mathbb{R}$, $\arg(G_{ij}) \in \{0, \pi\}$ --- only two values, with no intermediate states. There is no gauge-covariant derivative, no field strength, no forces. A universe over $\mathbb{R}$ has geometry but is dead: matter cannot interact. $\mathbb{R}$ is eliminated.

\medskip
The sole survivor is $\mathbb{C}$. Over $\mathbb{C}$:
\begin{itemize}
\item R1: $|G_{ij}| = |z|$ defines a norm. \checkmark
\item R2: $\arg(G_{ij}) \in (-\pi, \pi]$, and $\mathbb{C}^\times_1 = \text{U}(1) = S^1$ is a connected 1-dimensional Lie group. \checkmark
\item R3: $\mathbb{C}$ is commutative, so $\det(G)$ is well-defined and real-valued for Hermitian $G$. \checkmark \qedhere
\end{itemize}
\end{proof}

\subsection{Proof II: Topological confirmation}

The algebraic proof suffices, but R4 provides an independent topological confirmation.

\begin{theorem}[Uniqueness of $\mathbb{C}$ --- topological proof]
\label{thm:C_topology}
$\mathbb{C}$ is the unique normed division algebra whose phase group has $\pi_1 = \mathbb{Z}$.
\end{theorem}

\begin{proof}
The phase group (unit-norm elements) of each division algebra is a sphere:
\begin{equation}
\mathbb{R}: S^0, \qquad \mathbb{C}: S^1, \qquad \mathbb{H}: S^3, \qquad \mathbb{O}: S^7.
\end{equation}

The fundamental groups are:
\begin{align}
\pi_1(S^0) &= 0 &&\text{(disconnected: two isolated points, no loops)}, \\
\pi_1(S^1) &= \mathbb{Z} &&\text{(circle: loops wind an integer number of times)}, \\
\pi_1(S^3) &= 0 &&\text{(simply connected: all loops are contractible)}, \\
\pi_1(S^7) &= 0 &&\text{(simply connected)}.
\end{align}

Charge quantization requires that the holonomy around a closed loop takes discrete values. The winding number $Q = \frac{1}{2\pi}\oint d\phi$ must be an integer. This demands $\pi_1(\text{phase group}) \cong \mathbb{Z}$.

\begin{itemize}
\item \textbf{$\mathbb{R}$}: $\pi_1(S^0) = 0$. No loops exist (the space is disconnected). No gauge transport, no charges. Eliminated.
\item \textbf{$\mathbb{C}$}: $\pi_1(S^1) = \mathbb{Z}$. Loops wind integer times. The Dirac quantization condition $eg = n\hbar/2$ follows directly. Electric charge is quantized \emph{because} the phase group is a circle. \textbf{Unique survivor.}
\item \textbf{$\mathbb{H}$}: $\pi_1(S^3) = 0$. Every closed loop of phases is contractible to a point. No winding numbers, no charge quantization. (Note: $\pi_3(S^3) = \mathbb{Z}$ gives instantons, but these are non-perturbative tunneling events, not conserved particle charges.) Eliminated.
\item \textbf{$\mathbb{O}$}: $\pi_1(S^7) = 0$. Same argument. Also eliminated by non-associativity. \qedhere
\end{itemize}
\end{proof}

\subsection{The inner product structure}

Having established $\mathbb{K} = \mathbb{C}$, the relation $G_{ij}$ is most naturally an inner product:
\begin{equation}
G_{ij} = \langle\psi_i|\psi_j\rangle = \sum_{a=0}^{d-1} \psi^*_{i,a}\,\psi_{j,a}, \qquad \psi_i \in \mathbb{C}^d
\end{equation}
for some dimension $d$ to be determined. This form automatically satisfies:
\begin{enumerate}
\item Hermiticity: $G_{ji} = G^*_{ij}$ (so $G$ is a Hermitian matrix),
\item Positive semi-definiteness: $\mathbf{c}^\dagger G\,\mathbf{c} = \|\Psi^\dagger\mathbf{c}\|^2 \geq 0$,
\item Unit diagonal: $G_{ii} = \|\psi_i\|^2 = 1$ (normalization: no point is ``more real'' than another),
\item Rank constraint: $\text{rank}(G) \leq d$.
\end{enumerate}

The normalization $\|\psi_i\| = 1$ is not an additional assumption but a consequence of \emph{point democracy}: if $\|\psi_i\| \neq \|\psi_j\|$, then point $i$ would be ``more fundamental'' than point $j$, breaking the symmetry of the starting assumption.

\subsection{Physical interpretation}

$\mathbb{C}$ is the unique field with \emph{exactly one phase direction}. Not zero (that's $\mathbb{R}$: geometry without forces), not three (that's $\mathbb{H}$: forces without well-defined gravity), but one.

This single phase direction becomes $\text{U}(1)$ --- the electromagnetic gauge group --- the ``glue'' that will connect the spatial and temporal sectors of $\mathbb{C}^d$ (Sec.~\ref{ch:whyd5}). The existence of electromagnetism is not a contingent fact about our universe; it is a \emph{mathematical necessity} for any universe with both geometry and forces.

\begin{remark}
The hierarchy of division algebras mirrors the hierarchy of physical structures:
\begin{center}
\begin{tabular}{lcl}
$\mathbb{R}$ & $\to$ & Geometry only (general relativity without matter) \\
$\mathbb{C}$ & $\to$ & Geometry + forces + charges (the physical universe) \\
$\mathbb{H}$ & $\to$ & Forces without consistent gravity (inconsistent) \\
$\mathbb{O}$ & $\to$ & Non-associative (no matrix structure, no physics)
\end{tabular}
\end{center}
Only $\mathbb{C}$ sits in the ``Goldilocks zone'': rich enough for forces, constrained enough for gravity.
\end{remark}
